\begin{table}[H]
\centering\centering
\caption{Race and next-season roster retention within employer and season \label{tab:player-retention-employer-fe}}
\centering
\resizebox{\ifdim\width>\linewidth\linewidth\else\width\fi}{!}{
\fontsize{8}{10}\selectfont
\begin{tabular}[t]{lccc}
\toprule
  & (1) & (2) & (3)\\
\midrule
P(Black) & -0.032*** & -0.032*** & -0.028***\\
 & (0.008) & (0.008) & (0.008)\\
P(other race) & 0.002 & 0.002 & 0.006\\
 & (0.015) & (0.015) & (0.015)\\
\midrule
Observations & 46083 & 46082 & 46082\\
R$^2$ & 0.244 & 0.244 & 0.265\\
Within R$^2$ & 0.237 & 0.237 & 0.236\\
Position group $\times$ season FE & Yes & Yes & Yes\\
Employer $\times$ season FE & No & No & Yes\\
Career stage (A) & Yes & Yes & Yes\\
\shortstack[l]{Season-$t$ and career\\production (B, C)} & Yes & Yes & Yes\\
Pre-NFL signals (D, F) & Yes & Yes & Yes\\
Sample & \shortstack{All\\(Table \ref{tab:player-retention}, col. 4)} & \shortstack{Known\\employer} & \shortstack{Known\\employer}\\
Employer $\times$ season cells &  &  & 704\\
Mean of outcome & 0.778 & 0.778 & 0.778\\
Players & 11,298 & 11,298 & 11,298\\
P(Black) gap (pp) & -3.2 & -3.2 & -2.8\\
MDE (80\% power, pp) & 2.3 & 2.3 & 2.3\\
\bottomrule
\multicolumn{4}{l}{\rule{0pt}{1em}* p $<$ 0.1, ** p $<$ 0.05, *** p $<$ 0.01}\\
\end{tabular}}
\end{table}

{\footnotesize
\noindent\textit{Notes:} This table includes the estimation results of the column-(4) specification of Table \ref{tab:player-retention} (position group $\times$ season fixed effects and control blocks A, B, C, D and F) with, in column (3), employer $\times$ season fixed effects added, so that the race coefficients compare players of the same franchise in the same season. The outcome is 1 if the player is under contract to any franchise in at least one game week of season $t+1$ (retention anywhere in the league, not by the same franchise). The employer is the franchise of the player's last verified under-contract week of season $t$, measured from the season-$t$ weekly rosters only (never from $t+1$). When two franchises list a player under contract in that week (duplicate listings and trade weeks), the person-week deduplication picks one by a tie-break rather than by evidence; those player-seasons have no known employer and are excluded rather than assigned. Unverified 2016 preseason listings never set the employer. Sample flow: column (1) is the Table \ref{tab:player-retention} column (4) sample (46,083 player-seasons). Of these, 46,082 have a known employer (1 tied player-seasons dropped); 0 employer $\times$ season cells with a single player-season are dropped from columns (2) and (3) alike (0 player-seasons), because such a row is fitted exactly by the employer fixed effect in (3) and would carry weight in (2) only. Columns (2) and (3) are estimated on the identical 46,082 player-seasons (asserted on the estimation rows of both fits), 11,298 players, 32 employers and 704 employer $\times$ season cells. Columns (2) and (3) compare the race coefficient with and without employer fixed effects on fixed rows. This is specification sensitivity, not a causal decomposition of the gap: adding fixed effects also changes the residual race-score variation. The employer is the last recorded under-contract franchise during season $t$, not necessarily employment in the final week. Conditioning on that employer and observed season-$t$ employment leaves a descriptive comparison among selected players, not an employer treatment effect. The MDE row is the gap the column could detect with 80 percent power at the 5 percent level ($2.8$ times the clustered standard error), in percentage points. Employment is measured from the weekly rosters, one row per person-week (a player listed by two franchises in a week counts once), in game weeks of the listing franchise (bye weeks excluded, so counts are comparable across eras); under contract means on the 53-man roster, a reserve list or suspended/exempt, which excludes practice-squad, released, retired and free-agent listings. No listing at all in $t+1$ is an exit. The 2016 source leaks preseason rosters into weeks 1-2; listings of players who neither played that week nor stayed under contract from week 3 are unverified and do not count as employment, and a $t+1$ whose only listings are unverified, trade-pending or of unknown status leaves the outcome unknown (excluded) rather than recording an exit. Block A: experience-bin dummies, experience, age and its square. Block B: season-$t$ games played, injury weeks and position-appropriate production with position-group-specific slopes (as the lagged block of Table \ref{tab:pay-gap-learning}, measured in $t$). Block C: career production before $t$. Block D: log draft pick, undrafted, draft-round dummies, 247 rating and stars, combine measurables and athletic scores, college programme. Block F: CFBD pre-draft grade. Block E: season-$t$ snaps, snap-based starts and career snaps and starts (2013+; zero before with an indicator). Block G: weeks under contract in $t$, under contract in the final week of $t$ and more than one spell. Controls with incomplete coverage are zero-filled with a missing indicator, so no player-season is dropped for a missing control. Every column also includes fixed effects for the covariates of the race prediction's prior (position at NFL entry, entry era, draft-round bucket, college type and whether a home county is known), as regression calibration requires. Under the predicted measure the coefficient on P(Black) equals the Black-white gap only under two assumptions beyond names and hometown being unrelated to the outcome given race and the controls: (i) P(Black) is calibrated conditional on every control in the column, not only on the prior's covariates, and (ii) the gap is the same across values of the controls (otherwise the coefficient is a variance-weighted average of gaps). Adding controls changes which conditional calibration is required; it also changes the overlap of P(Black) within cells, which affects precision rather than the estimand. Changes in the coefficient across columns therefore cannot be attributed to the controls' effect on the gap without checking (i) and (ii) in each column. The estimates are conditional associations among players already employed in season $t$; survival to $t$ is itself selected on the same unobservables that drive the outcome, so the coefficients describe who stays and who reaches a bargained contract given observed quality, not a causal effect of race. Pay is neither an outcome nor a control; roster weeks are not paid weeks. The P(Black) gap is $100\hat\beta_1$, in percentage points of the outcome. Standard errors, clustered by player, in parentheses. $^{*}p<0.1$, $^{**}p<0.05$, $^{***}p<0.01$. P(Black) and P(other race) enter together, so the coefficient on P(Black) compares a Black (predicted: non-Hispanic Black alone) with a white player. Race is predicted, not observed: each person's probability of being non-Hispanic Black alone combines first name, surname and hometown county (BIFSG) with an NFL-specific prior estimated by EM on predetermined characteristics (players: position at entry, rookie era, draft round, college type, county availability; staff: role, unit and era at first appearance); documented race statements are not used (notes/race-prediction-design.md). Black Hispanic and multiracial Black persons are in the other-race probability, so P(Black) is narrower than the Black-alone-or-in-combination definition of the hand codes and of TIDES. Regressions use the probabilities (regression calibration) and control for the prior's covariates; the coefficient on P(Black) is the Black-white gap if the probabilities are calibrated given the regression's controls and names and hometown are unrelated to the outcome given race and the controls. The validation exhibit compares the predictions with published TIDES shares.
\par}

